\originalchapter{核酸的解链与二级结构}
\originalsection{DNA中的环熵}
双链DNA的每段配对有结合权重 $w^\ell$，熔解气泡由两条单链连接同样的两个端点。对理想三维随机链，返回原点的概率按 $\ell^{-3/2}$ 衰减；更一般的环权重写成 $B(\ell)=A s^\ell/\ell^c$。指数 $c$ 包含环内和环外的排除体积效应，不能把三维高斯值用于所有DNA模型。

令总长度的共轭变量为 $z$，有 $\widetilde R(z)=\sum_{\ell\ge1}(wz)^\ell=wz/(1-wz)$，$\widetilde B(z)=A\sum_{\ell\ge1}(sz)^\ell/\ell^c$。双链段与气泡交替排列，几何级数给出的主体分母是 $1-\widetilde R\widetilde B$；两端的规定只改变分子。自由能密度由最靠近原点的正实奇点 $z_*$ 决定，$f=\kb T\ln z_*$。

低温时奇点是 $\widetilde R\widetilde B=1$ 的极点，高温时由气泡级数的收敛边界 $z=1/s$ 控制。若 $c\le1$，环和在边界发散，极点不在有限参数下撞到边界；若 $1<c\le2$，环和有限但平均环长发散，配对比例连续趋零；若 $c>2$，平均环长有限，比例可以跳变。这一分类来自 $\sum\ell^{-c}$ 与 $\sum\ell^{1-c}$ 的收敛性质，而非把一维系统简单等同于有无相变。

\originalsection{RNA配对递推}
单链RNA允许内部配对。在不考虑假结的模型中，一段配对不能与另一段交叉。这不是全部真实RNA结构，而是一类可递推求和的二级结构。令 $Z_{ij}$ 为从 $i$ 到 $j$ 的分段配分函数，末端 $j$ 或未配对，或与 $k$ 配对。若配对权重为 $q_{kj}$，最小环长为 $h$，则
\[
Z_{ij}=Z_{i,j-1}+\sum_{k=i}^{j-h-1}Z_{i,k-1}q_{kj}Z_{k+1,j-1},\qquad Z_{i,i-1}=1.
\]
因为禁止交叉，$k,j$ 内部和左侧区间可以独立求和。递推输出全构象统计权重，和只求最低能量的算法不同。真实能量还包括堆叠、环长、盐浓度等参数；本式的简单权重适合解释算法结构，不能直接当作实际RNA预测器。

\originalsection{拉力与核酸解链}
力控制条件下，应使用 $G=H-FX$，而固定延伸时应使用 Helmholtz 自由能。两种系综中的稳定条件相互关联，但跳变区间的观测方式不同。DNA的气泡可以承受某些拉伸几何下的两链载荷；RNA内部配对团块则常被建模为不承受骨架拉力。这使两种解链的配分函数出现相似的交替结构，却有不同的力权重。

设不配对骨架段 $P(\ell)=A^{-1}g^\ell\exp[\beta^2F^2a^2\ell/6]$，配对团块 $R(\ell)=f^\ell/\ell^{3/2}$。在 $z_c=1/f$ 处，团块和有限但导数发散。当极点撞到此边界时，开始连续解链。所选原题中给出临界力、线性起始奇性和此高斯链近似的适用范围。DNA受力完整相图与旧在线RNA服务任务没有纳入本版逐题解答。
